% Corrected French for printed p.168; departures are append-only below.
\SourcePage{168}

\LayerTarget{tohoku:proof:3.6.1:p2}
\noindent de sous-faisceaux $F_j$ dont chacun admet une famille \textit{finie}
de générateurs.

\LayerTarget{tohoku:construction:finite-generator-reduction}
Comme $C''$ est épaisse et stable par sommes directes infinies, elle est aussi
stable par limites inductives (car une limite inductive d'objets dans $C'$ est
isomorphe à un quotient de leur somme directe) ; comme
\[
T^p(F)=\varinjlim_j T^p(F_j),
\]
il suffit, pour prouver $T^p(F)\in C''$, de prouver que
$T^p(F_j)\in C''$ pour tout $j$, ce qui nous ramène au cas où $F$ admet une
famille \textit{finie} de générateurs $(f_i)_{1\leq i\leq k}$.

\LayerTarget{tohoku:construction:finite-generator-filtration}
Désignons alors par $F_n$ ($0\leq n\leq k$) le sous-faisceau de $F$ engendré
par les $f_i$ avec $1\leq i\leq n$. Les $F_n$ forment une suite croissante
finie de sous-faisceaux de $F$ dont les quotients successifs
$F_n/F_{n-1}$ ($1\leq n\leq k$) admettent chacun \textit{un} générateur.

\LayerTarget{tohoku:eq:finite-generator-quotient-reduction}
Raisonnant par récurrence sur la longueur $k$ de cette suite, et utilisant le
fait que $T^p$ est semi-exact, on est ramené à prouver que
$T^p(F_n/F_{n-1})\in C''$ pour tout $n$.

\LayerTarget{tohoku:eq:single-generator-presentation}
Cela nous ramène au cas où $F$ est engendré par \textit{un} générateur,
c'est-à-dire où il existe une suite exacte de faisceaux
\[
0\longrightarrow R\longrightarrow Z_U\longrightarrow F\longrightarrow0.
\]

\LayerTarget{tohoku:eq:single-generator-cohomology}
où $R$ est un sous-faisceau de $Z_U$ donc aussi de $Z$. On en conclut la suite
exacte
\[
T^p(Z_U)\longrightarrow T^p(F)\longrightarrow T^{p+1}(R),
\]
et comme on a $T^p(Z_U)\in C''$ par hypothèse, on est ramené à prouver que
$T^{p+1}(R)\in C''$.

\LayerTarget{tohoku:construction:constant-section-generators}
Or le sous-faisceau $R$ du faisceau constant $Z$ est engendré par une famille
$(f_i)_{i\in I}$ de sections \textit{constantes} $n_i$ de $Z$ sur des ouverts
$U_i$. Procédant comme plus haut, on est ramené au cas où cette famille est
finie. On peut bien entendu supposer les $n_i>0$.

\LayerTarget{tohoku:construction:gcd-refinement}
Nous pouvons supposer de plus les $f_i$ choisis de telle façon que pour tout
$x\in X$, il existe parmi les $n_i$ pour lesquels $U_i$ contient $x$, un
générateur du sous-groupe $R(x)$ du groupe $Z(x)={}$ groupe $Z$ des entiers.
Il suffit pour ceci, pour toute partie $(i_1,\ldots,i_p)$ de l'ensemble $[1,k]$
des $k$ premiers entiers $>0$, de considérer la section
$f_{i_1\cdots i_p}$ de $R$ sur $U_{i_1}\cap\cdots\cap U_{i_p}$ dont la valeur
est le $\operatorname{pgcd}$ de $n_{i_1},\ldots,n_{i_p}$ et d'adjoindre ces
sections au système de générateurs.

\LayerTarget{tohoku:construction:R-filtration}
Supposons donc notre hypothèse sur les $f_i$ satisfaite, et supposons
$n_1\leq n_2\leq\cdots\leq n_k$. Si on désigne par $R_m$ ($0\leq m\leq k$)
le sous-faisceau de $R$ engendré par les sections $f_i$ avec
$1\leq i\leq m$, les $R_m$ forment une suite croissante finie de
sous-faisceaux de $R$.

\LayerTarget{tohoku:eq:R-filtration-reduction}
Et on montre encore comme plus haut, par récurrence sur $k$, et utilisant le
fait que $T^{p+1}$ est semi-exact, que pour prouver $T^{p+1}(R)\in C''$, il
suffit de prouver $T^{p+1}(R_m/R_{m-1})\in C''$ pour $1\leq m\leq k$.

\LayerTarget{tohoku:construction:Vm-Ym}
Or soit pour tout $m$, $V_m$ la réunion des $U_i$ pour $1\leq i\leq m$, et
soit
\[
Y_m=U_m\cap \mathrm C V_{m-1}.
\]

\LayerTarget{tohoku:eq:R-quotient-ZYm}
Je dis que $R_m/R_{m-1}$ est \textit{isomorphe à} $Z_{Y_m}$.

\LayerTarget{tohoku:proof:R-quotient-ZYm:p1}
En effet sa restriction au complémentaire de $V_m$ est évidemment nulle
(puisque celle de $R_m$ l'est déjà), il en est de même de sa restriction à
$V_{m-1}$. Il suffit de le vérifier en tout $x\in V_{m-1}\cap U_m$, or par
hypothèse, parmi les $n_i$ relatifs aux $U_i$ contenant $x$ figure leur
$\operatorname{pgcd}$ $n_{i_0}$, qui doit donc diviser à la fois $n_m$ et un
$n_i$ avec $i<m$ (prendre un indice $i<m$ avec $x\in U_i$, qui existe puisque
$x\in V_{m-1}$). Alors ou bien $i_0<m$, ce qui prouve que
$f_m(x)\in R_{m-1}(x)$, ou $i_0\geq m$, d'où $n_{i_0}\geq n_m$ et par suite
$n_{i_0}=n_m$, donc $n_m$ (divisant $n_i$ et $\geq n_i$) est égal à $n_i$,
d'où encore $f_m(x)\in R_{m-1}(x)$, donc dans les deux cas
$R_{m-1}(x)=R_m(x)$. Ainsi la restriction de $R_m/R_{m-1}$ à la réunion de
$\mathrm C V_m$ et de $V_{m-1}$, i.\,e. à $\mathrm C Y_m$, est

\TohokuContinuation{169}
